\section{AEP、典型集与信源编码定理}
\label{sec:07-Information-Theory/04-Source-Coding/04}

前两节都在往下界靠：Huffman 贴到熵加一比特，算术编码贴到熵加两比特除以 \(n\)。逼近的一侧已经交代清楚了。可下界本身目前只有一个逐符号的论证——它说的是``用分布 \(Q\) 编码 \(P\) 要多付 \(\KL(P\|Q)\)''，对单个符号成立，但没有回答一个更朴素的问题：把一千个符号当成一个整体，凭什么它就装不进少于 \(1000H\) 个比特？

这一节换一种数法。不算长度，数个数。

\subsection{长序列的自信息是一个样本均值}

长度为 \(n\) 的独立同分布序列，概率是各项相乘：

\[
p(X^n)=\prod_{i=1}^{n}p(X_i).
\]

取负对数再除以 \(n\)，乘积变成算术平均：

\[
-\frac1n\log p(X^n)=\frac1n\sum_{i=1}^{n}\bigl[-\log p(X_i)\bigr].
\]

右边每一项 \(-\log p(X_i)\) 都是随机变量，期望恰是 \(H(X)\)。于是这个量就是一列独立同分布随机变量的样本均值，大数定律直接接管：

\[
-\frac1n\log p(X^n)\longrightarrow H(X)\quad\text{几乎必然}.
\]

换成概率的形式，绝大多数序列满足 \(p(x^n)\approx 2^{-nH(X)}\)。这叫渐近等分性质（Asymptotic Equipartition Property, AEP）。``等分''说的不是所有序列概率相等——全零序列和一半零一半一的序列概率相差极远——而是那些实际会出现的序列，在一阶指数尺度上概率同量级，差异被指数抹平了。

\subsection{典型集把极限定理变成可以数的东西}

给定 \(\varepsilon>0\)，把``接近''固定下来，定义典型集

\[
A_\varepsilon^{(n)}=\set{x^n:\Bigl|-\tfrac1n\log p(x^n)-H(X)\Bigr|<\varepsilon}.
\]

AEP 等价于 \(P\bigl(X^n\in A_\varepsilon^{(n)}\bigr)\to1\)。典型集里每个序列的概率被夹在 \(2^{-n(H+\varepsilon)}\) 与 \(2^{-n(H-\varepsilon)}\) 之间，这两个夹子把集合的大小也卡住了：总概率不超过一，所以元素数至多约 \(2^{n(H+\varepsilon)}\)；集合又承载着接近一的概率，而每个元素概率至多 \(2^{-n(H-\varepsilon)}\)，所以元素数至少约 \(2^{n(H-\varepsilon)}\)。让 \(\varepsilon\) 趋于零，在一阶指数率上

\[
\bigl|A_\varepsilon^{(n)}\bigr|\doteq 2^{nH(X)}.
\]

记号 \(\doteq\) 表示两边取对数除以 \(n\) 后极限相同，有限 \(n\) 时还浮动着 \(2^{\pm n\varepsilon}\) 的因子。

关键在于这两个指数的对比：全部序列有 \(2^{n\log m}\) 条，典型的只有 \(2^{nH}\) 条，而 \(H<\log m\) 只要分布不均匀就成立。

\begin{center}
\begin{tikzpicture}[x=1cm,y=1cm,font=\small,>=stealth]
  \draw[black!50,fill=black!4] (0,0) circle (2.1);
  \fill[blue!35] (0,0) circle (0.34);
  \draw[black!55] (0,0) circle (0.34);
  \node[right,font=\footnotesize] at (0.45,0.05) {典型集 \(\approx 2^{nH}\) 条};
  \node[font=\footnotesize,black!60] at (0,-1.55) {全部 \(2^{n}\) 条序列};
  \node[above,font=\footnotesize] at (0,2.2) {按条数看：几乎可以忽略};
  \draw[black!50,fill=black!4] (6.6,0) circle (2.1);
  \fill[blue!35] (6.6,0) circle (2.0);
  \draw[black!55] (6.6,0) circle (2.0);
  \node[font=\footnotesize] at (6.6,0.05) {概率质量 \(\to1\)};
  \node[above,font=\footnotesize] at (6.6,2.2) {按概率看：几乎就是全部};
  \draw[->,black!55] (9.55,1.62) -- (8.06,1.46);
  \node[right,font=\footnotesize,black!60,align=left] at (9.6,1.62) {其余序列合起来\\概率 \(\to0\)};
\end{tikzpicture}

\emph{同一个典型集，两把尺子量出完全相反的印象。压缩就是利用这个反差：只给左图那一小块编号，右图告诉你这样做几乎不会出错。}
\end{center}

拿 Bernoulli\((0.99)\)、\(n=1000\) 代进去感受一下量级。全部序列 \(2^{1000}\) 条，典型的约 \(2^{81}\) 条。给这 \(2^{81}\) 条各配一个八十一位的索引，再留一个特殊码字统一处理所有非典型序列（它们加起来的概率趋于零，碰上就报错），每符号平均花掉 \(0.081\) 比特——逐符号编码那个下不去的一比特，一下子甩开了十倍。

\subsection{两个方向的论证都只需要数个数}

可达性就是刚才那个构造。对任意速率 \(R>H(X)\)，取 \(\varepsilon<R-H\)，块长足够大时典型集不超过 \(2^{nR}\) 条，索引装得下，而出错概率就是落在典型集外的概率，趋于零。

不可达的一侧同样是计数。若速率 \(R<H(X)\)，长度 \(nR\) 的索引最多区分 \(2^{nR}\) 条序列。就算精挑细选最有利的那些典型序列来编码，它们能覆盖的概率至多是

\[
2^{nR}\cdot 2^{-n(H-\varepsilon)}=2^{-n(H-R-\varepsilon)},
\]

取 \(\varepsilon<H-R\) 后这个指数趋于零；非典型部分本来概率也趋于零。两块加起来，成功恢复的概率趋于零。

\begin{theorem}
离散无记忆信源 \(X\) 的近无损压缩存在一个临界速率，恰为 \(H(X)\)：任何 \(R>H(X)\) 都可达，且错误概率随块长趋于零；任何 \(R<H(X)\) 不可达。
\end{theorem}

定理的价值不在于它给出一个数，而在于它把``这份数据还能不能再压''从一个经验问题变成了有确定答案的问题。以前只能说``我试了几种压缩器，最好的到这个大小''；有了这条定理，只要模型确定，就能说``低于这个数不可能，而且高一点点一定可以''。

条件也要看清。定理成立于一个明确的渐近设定：信源模型已知，块长趋于无穷，允许错误概率趋于零。它没有承诺任何有限文件能恰好压到 \(nH(X)\) 比特，更没有承诺任意数据都可压缩。真实压缩器还要额外支付三笔：模型未知得先学，模型描述本身占比特，有限块长带来的离散损失。

从算术编码的角度看，这段论证会更有实感。典型序列的概率约 \(2^{-nH}\)，它在 \([0,1)\) 上占的区间宽度也约 \(2^{-nH}\)，用二进制定位这个区间恰需约 \(nH\) 位。典型集编码和算术编码在指数阶上是同一件事，差别只在常数与实现。

\subsection{典型不是``最可能''}

有个坑值得单独标出来：典型集不是概率最大的那些序列的集合。对 Bernoulli\((0.99)\)，全零序列通常是单条概率最高的序列，但它往往不典型——它的经验频率是 \(100\%\) 的零，离真实比例 \(99\%\) 有距离，自信息因此偏离 \(H(X)\)。典型集关心的是整体：每个成员未必名列前茅，但整个集合承载了几乎全部概率。

公平硬币是另一端的边界。每条长度 \(n\) 的序列概率都是 \(2^{-n}\)，熵是一比特，于是所有序列都典型，没有低概率区域可以丢弃。AEP 在这里没有提供任何压缩杠杆，而这正是对的——最大熵的信源本来就没有冗余可挤。

\subsection{有记忆的信源换一个熵}

上面每一步都用了独立同分布。若信源有依赖，边缘熵 \(H(X_1)\) 不再是压缩极限，正确的量是熵率

\[
\bar H=\lim_{n\to\infty}\frac1n H(X_1,\dots,X_n).
\]

在平稳遍历的条件下，Shannon--McMillan--Breiman 定理给出对应的 AEP：

\[
-\frac1n\log p(X_1,\dots,X_n)\longrightarrow\bar H\quad\text{几乎必然}.
\]

典型集的大小随之变成约 \(2^{n\bar H}\)，编码定理把 \(H(X)\) 换成 \(\bar H\) 后一字不改地成立。自然语言的熵率远低于单字符熵，这正是语言模型能把每字符比特数压到远低于逐字符频率所允许的水平的原因——多出来的压缩全部来自上下文依赖。

若信源既非平稳也非遍历，单一典型集的图景可能要换成若干典型集的混合，甚至失去意义。AEP 是独立同分布假设送的礼物，在更野的信源面前，概率集中在哪里得自己去找。

至此无损压缩的账算完了：熵是底线，底线可以逼近，底线之下什么也做不到。但整套论证有一个从未松动的前提——解码必须逐比特还原原序列。照片、音频、视频从来不需要这种精确度，人眼看不出的差别没有理由付比特去保留。下一节把这个前提松开，看看下界会掉到哪里。
